%
% Slides for given talks:
%
% * 2007-02-26: PyCon 2007, Dallas (Texas), USA
%
%\documentclass[compress,trans]{beamer}
\documentclass{beamer}

\usepackage{graphicx}

\mode<presentation>
{
\usetheme{Singapore} % Singapore, default
  \setbeamertemplate{frametitle}[default][left]
  \setbeamercovered{transparent}
}

\usepackage[english]{babel}
\usepackage[latin1]{inputenc}

\usepackage{times}
\usepackage[T1]{fontenc}

\title{An extension to DBAPI 2.0 \\
for easier SQL queries}
\subtitle{}

\author{Martin Blais}

\institute{}

\date{PyCon 2007}

\subject{dbapiext presentation @ PyCon 2007}

\setcounter{tocdepth}{4}

\newcommand{\todo}[1]{}


\AtBeginSubsection[]
{
  \begin{frame}<beamer>
    \frametitle{Outline}
    \tableofcontents[currentsection,currentsubsection]
  \end{frame}
}


\begin{document}

%-------------------------------------------------------------------------------
\begin{frame}
  \titlepage
\end{frame}

%===============================================================================


%-------------------------------------------------------------------------------
\begin{frame}[fragile]
  \frametitle{Introduction}

  DBAPI's \texttt{Cursor.execute()} method interface is inconvenient to use.

\vfill

  General form:
\begin{verbatim}
  cursor.execute(<string>, <tuple-or-dict>)
\end{verbatim}

\vfill

\end{frame}



Raymond's comments:


  show substitution results
  start from the POV of the escaped substitution first, work your way back
  make sure it is worded next to the examples why they fail

  then, show how some pieces I don't want to escape, have to be used with %%s

  introduce keyword arguments separately

  show an example for the problem with lists

    my solution: is to accept a list argument


questions:

- "query writing" unclear
- prepared statements
- could support other %S specifiers
- use cheetah templating

indent Mickey mouse query (split it, ray)

- add a little graphics to show that it only sits over dbapi

people are not familiar with the problem, start with the erroneous statements
examples, and work your way to the ugly execute() method.

then present the improvemetns.












%-------------------------------------------------------------------------------
\begin{frame}[fragile]
  \frametitle{\texttt{Cursor.execute()} interface}


\vfill
\pause









  Common mistake \#1:
\begin{verbatim}
 cursor.execute(
   "INSERT INTO Users (username) VALUES (%s)" % 
   var_username)
\end{verbatim}


%%   For example:

%% \begin{verbatim}
%%     INSERT INTO Users (username, email) 
%%        VALUES ('George', 'g@soros.com');
%% \end{verbatim}

%% \begin{verbatim}
%%   cursor.execute("""
%%     INSERT INTO %s (%s, %s) VALUES (%%s, %%s)
%%     """ % ("Users", "username", "email"),
%%     (var_username, var_email))
%% \end{verbatim}

\end{frame}


%-------------------------------------------------------------------------------
\begin{frame}[fragile]
  \frametitle{\texttt{Cursor.execute()} interface}

%%   These are \emph{also} incorrect.

  Common mistake \#2:
\begin{verbatim}
 cursor.execute(
   "INSERT INTO Users (username) VALUES ('%s')" % 
   var_username)
\end{verbatim}

\pause

  Common mistake \#3:
\begin{verbatim}
 cursor.execute(
   "INSERT INTO Users (username) VALUES (%s)" % 
   repr(var_username))
\end{verbatim}

\end{frame}

  Common mistake \#4:
\begin{verbatim}
 cursor.execute(
   "INSERT INTO Users (username) VALUES (%s)",
   var_username)
\end{verbatim}


%-------------------------------------------------------------------------------
\begin{frame}[fragile]
  \frametitle{\texttt{Cursor.execute()} interface}

  You \emph{must} let DBAPI do its database-specific escaping of values:
\begin{verbatim}
 cursor.execute(
   "INSERT INTO Users (username) VALUES (%s)",
   (var_username,))
\end{verbatim}

String constants, timestamps, dates, etc. ; Formats vary depending on the
database.

\pause
\vfill
Problems:
\begin{itemize}
\item Two lists of parameters is error-prone
\item You have to provide a tuple or a dict for the argument
\item It does not understand lists 
\item You can't leverage the power of keyword arguments
\end{itemize}

\end{frame}


%-------------------------------------------------------------------------------
\begin{frame}[fragile]
  \frametitle{Introduction}

\vfill
  With this work:
  \begin{itemize}
  \item Provide a simple extension that gets rid of the pitfalls
  \item Make it much easier to write queries
  \item A single pure Python module
  \item Support a number of DBAPI implementations
  \item Deals only with query writing
  \end{itemize}

\end{frame}



%-------------------------------------------------------------------------------
\begin{frame}[fragile]
  \frametitle{\texttt{Cursor.execute()} interface}

  When you write real-world queries (instead of Mickey-mouse example
  queries), it gets even messier:
\begin{verbatim}
  cursor.execute("""
    SELECT %s FROM %s 
      WHERE %s > %%s LIMIT %s
    """ % (','.join(columns), "Users", 
           "age", 10),
    (18,))
\end{verbatim}

\begin{itemize}
\item Because of string interpolation, you have to \textbf{double-escape} the
  format specifiers for the escaped values
\item The parameters in the strings are in a \textbf{different order} than the
  function arguments (easy to make mistakes!)
\end{itemize}

\end{frame}




%-------------------------------------------------------------------------------
\begin{frame}[fragile]
  \frametitle{New format specifier (\%S)}

  We provide a new \texttt{execute()} method, which supports a format specifier
  for escaped arguments: \texttt{\%S} (capital S)

\begin{verbatim}
 cursor.execute_f(
   "INSERT INTO Users (username) VALUES (%S)",
   var_username)
\end{verbatim}

  You can now mix vanilla and escaped values in the arguments:
\begin{verbatim}
 cursor.execute_f(
   "INSERT INTO Users (%s) VALUES (%S)",
   "username", var_username)
\end{verbatim}

\end{frame}



%-------------------------------------------------------------------------------
\begin{frame}[fragile]
  \frametitle{Lists are understood}

  Lists are automatically joined with commas:
\begin{verbatim}
 columns = ["username", "email", "age"]
 cursor.execute_f(
   "INSERT INTO Users (%s) VALUES (...)",
   columns, ...)
\end{verbatim}

\vfill
\begin{verbatim}
   INSERT INTO Users 
      ('username', 'email', 'age') 
      VALUES (...)
\end{verbatim}
\vfill

\end{frame}



%-------------------------------------------------------------------------------
\begin{frame}[fragile]
  \frametitle{Lists are understood}

  This also works for escaped arguments:
\begin{verbatim}
 columns = ["username", "email", "age"]
 values = [var_username, var_email, var_age]
 cursor.execute_f(
   "INSERT INTO Users (%s) VALUES (%S)",
   columns, values)
\end{verbatim}

\vfill
\begin{verbatim}
   INSERT INTO Users 
      ('username', 'email', 'age') 
      VALUES ('Warren', 'w@buffet.com', 76)
\end{verbatim}
\vfill

\begin{itemize}
\item Values are escaped individually and then comma-joined
\end{itemize}

\end{frame}


%-------------------------------------------------------------------------------
\begin{frame}[fragile]
  \frametitle{Dictionaries are understood}

  Dictionaries are rendered as required for \texttt{UPDATE} statements:
  \begin{itemize}
  \item Comma-separated \texttt{<name> = <value>} pairs
  \item Values are DB-escaped automatically
  \end{itemize}

\begin{verbatim}
  UPDATE languages 
    SET id = 3, brazil = 'portuguese'
\end{verbatim}
\vfill

\begin{verbatim}
  values = {"id": 3, 
            "brazil": "portuguese"}

  cursor.execute_f("UPDATE languages SET %S", 
                   values)
\end{verbatim}
\vfill


(Suggestion by D. Mertz)

\end{frame}



%-------------------------------------------------------------------------------
\begin{frame}[fragile]
  \frametitle{Positional and Keywords Arguments}

  Positional and keyword arguments can be used simultaneously:

\begin{verbatim}
  cursor.execute_f("""

     SELECT %(table)s FROM %s 
       WHERE id = %(id)S

  """, column_names, table=tablename, id=42)
\end{verbatim}

\begin{itemize}
\item You can recycle arguments this way \\
(i.e. a table or column name that occurs multiple times)
\end{itemize}


\end{frame}



%-------------------------------------------------------------------------------
\begin{frame}[fragile]
  \frametitle{Performance and Remarks}

  \begin{itemize}
  \item The extension only massages your query in a form that can be digested by
    DBAPI's \texttt{Cursor.execute()}

  \item I lied slightly in my examples, you have to use it like this:
\begin{verbatim}
   execute_f(cursor, """
      ...
\end{verbatim}

  \item We cache as much of the preprocessing as possible \\
(similar to \texttt{re}, \texttt{struct})
    \begin{itemize}
    \item You can cache your queries at load time with \texttt{qcompile()}.
    \end{itemize}
  \end{itemize}


\end{frame}



%-------------------------------------------------------------------------------
\begin{frame}[fragile]
  \frametitle{Future work}

Ideally, we would want to automatically parse the SQL queries and
determine which arguments should be quoted 

  \begin{itemize}
  \item A lot more work
  \item Would have to be done at load time for performance reasons
  \end{itemize}
  

\end{frame}



%-------------------------------------------------------------------------------
\begin{frame}[fragile]
  \frametitle{Questions}

  \begin{center}


{\Large
\texttt{dbapiext} is part of a \\
package named \texttt{antiorm}
}

\vfill

{\LARGE
antiorm homepage: \\
\verb=http://furius.ca/antiorm/=
}

\vfill

{\LARGE Questions?}

  \end{center}

\end{frame}


%===============================================================================
\end{document}




%% %-------------------------------------------------------------------------------
%% \begin{frame}[fragile]
%%   \frametitle{Motivation (Anti-ORM Activism)}
%% 
%%   \begin{center}
%%     \includegraphics[width=0.8\textwidth]{noorms.png}
%%   \end{center}
%% \end{frame}
%% 

%% %-------------------------------------------------------------------------------
%% \begin{frame}[fragile]
%%   \frametitle{Motivation}
%% 
%%   \begin{quote}
%%     While some like to use ORMs to interface with databases, we prefer to
%%     maintain the full power of the SQL language, and instead focus on making the
%%     building of SQL queries much more convenient.
%%   \end{quote}
%% 
%%   \begin{itemize}
%%   \item The domain-specific mini-language of SQL can never be fully replaced
%%     with an API
%%   \item ORMs hide the important facts: when does data transfer occur?  Does data
%%     get queried multiple times?
%%   \end{itemize}
%% 
%% \end{frame}
%% 
